\documentclass[10pt]{ctexart}
\usepackage[a4paper,margin=1in]{geometry}
\usepackage{amsmath,amssymb,mathtools,bm}
\usepackage{booktabs,longtable,array}
\usepackage{enumitem}
\usepackage{listings}
\usepackage{xcolor}
\usepackage{microtype}
\usepackage{hyperref}
\setlist[itemize]{leftmargin=1.5em}
\setlist[enumerate]{leftmargin=1.8em}
\lstset{basicstyle=\ttfamily\small,breaklines=true,columns=fullflexible,keepspaces=true,frame=single,framerule=0.2pt,xleftmargin=0.5em,xrightmargin=0.5em}
\newcommand{\clip}{\operatorname{clip}}
\newcommand{\uid}{\operatorname{uid}}
\newcommand{\trajid}{\operatorname{traj\_uid}}
\newcommand{\stepidx}{\operatorname{step\_idx}}
\title{013：当前项目进展与 Value 收敛分析}
\author{}
\date{2026-06-09}
\begin{document}
\maketitle
\vspace{-1.2em}

本文记录当前 \texttt{verl-agent} 项目的实验进展，并重点分析 StepPPO / GiGPO value-aux 方向中 value head 的收敛情况。分析基于当前仓库内的脚本、文档和 console log，不引入外部结果。

\section*{1. 当前项目进展}

当前项目主线已经从单纯复现 GiGPO baseline，推进到 StepPPO 及 shared value head 的稳定性诊断阶段。

\subsection*{1.1 已完成或已有有效日志的实验}

\begin{small}
\begin{longtable}{@{}p{0.23\linewidth}p{0.23\linewidth}p{0.23\linewidth}p{0.23\linewidth}@{}}
\toprule
实验 & 日志 & 状态 & 关键结论 \\
\midrule
GiGPO seed 2026 & \texttt{logs/webshop\_gigpo\_qwen25\_15b\_4gpu\_paper\_align\_simple\_20260530\_114442.log} & 完整 250 step & 强基线，最终 task/success/text = 0.903/0.812/7.963 \\
GiGPO seed 2077 & \texttt{logs/webshop\_gigpo\_qwen25\_15b\_4gpu\_paper\_align\_simple\_20260531\_004613.log} & 完整 250 step & 最终 task/success/text = 0.883/0.754/7.341 \\
GiGPO seed 2501 & \texttt{logs/webshop\_gigpo\_qwen25\_15b\_4gpu\_paper\_align\_simple\_20260531\_123134.log} & 完整 250 step & 最终 task/success/text = 0.869/0.691/6.952 \\
StepPPO-v1 step\_norm seed 2026 & \texttt{logs/webshop\_step\_ppo\_qwen25\_15b\_4gpu\_paper\_align\_simple\_20260604\_102927.log} & 跑到 step 187 / val 185 & 明显弱于 GiGPO，value loss 较大，v1 不稳定性已记录于 \texttt{EXPS/009\_step\_ppo\_instability\_analysis.md} \\
StepPPO-v2 seed 2026 & \texttt{logs/webshop\_step\_ppo\_qwen25\_15b\_4gpu\_paper\_align\_simple\_20260605\_031621.log} & 完整 250 step & 相比 v1 后期 value 有所缓和，但 valid action ratio 低，validation 仍未追上 GiGPO \\
GiGPO-v1 value-aux detach=false seed 2026 & \texttt{logs/webshop\_gigpo\_v1\_value\_aux\_detach\_false\_qwen25\_15b\_4gpu\_paper\_align\_simple\_20260608\_115519.log} & 跑到 step 204 / val 200 & policy advantage 保持 GiGPO，但 auxiliary value loss 仍带来明显退化和局部崩盘 \\
\bottomrule
\end{longtable}
\end{small}


三个 GiGPO baseline 已形成较稳定参照：step 250 的平均 task score 约 0.885，平均 success rate 约 0.752，平均 text score 约 7.419。StepPPO-v1/v2 和 GiGPO-value-aux 目前都没有达到该水平。

\subsection*{1.2 当前算法与脚本状态}

\begin{itemize}
\item StepPPO-v1 使用 \texttt{normalize\_episode\_advantage=True}、\texttt{normalize\_step\_advantage=True}、\texttt{step\_advantage\_w=1.0}、\texttt{value\_loss\_coef=0.1}，且 \texttt{detach\_value\_backbone=False}。
\item StepPPO-v2 实际脚本为 \texttt{exps/run\_webshop\_step\_ppo\_v2\_4gpu\_paper\_align\_simple.sh}，核心配置是 \texttt{step\_advantage\_w=0.5}、\texttt{episode\_mode=mean\_norm}、\texttt{final\_advantage\_mode=direct}、\texttt{normalize\_episode\_advantage=False}、\texttt{normalize\_step\_advantage=True}、\texttt{normalize\_final\_advantage=False}、\texttt{value\_loss\_coef=0.03}、\texttt{detach\_value\_backbone=False}、\texttt{ppo\_micro\_batch\_size\_per\_gpu=8}。
\item GiGPO-v1 value-aux 由 \texttt{verl/trainer/ppo/gigpo\_v1\_value\_aux\_trainer.py} 接入：policy advantage 仍是原 GiGPO，额外把 StepPPO-style \texttt{step\_returns} 塞给 shared value head 做 auxiliary regression。
\item \texttt{exps/run\_webshop\_gigpo\_value\_aux\_ablation\_grpo\_seed2026.sh} 已规划 \texttt{detach\_false -> detach\_true -> GRPO} 的 batch ablation，但当前有效 full/long 日志主要是 \texttt{detach\_false}；没有看到正式完整 4GPU GRPO baseline 日志。
\item \texttt{EXPS/tex/011\_step\_ppo\_v2\_improvements.tex} 中早期叙述曾提到 detach value backbone；但以当前脚本、日志和 \texttt{EXPS/tex/012\_step\_ppo\_v2\_paper\_method.tex} 为准，当前 v2 是 \texttt{detach\_value\_backbone=False}。
\end{itemize}

\section*{2. Validation 进展概览}

\begin{small}
\begin{longtable}{@{}p{0.18\linewidth}p{0.18\linewidth}p{0.18\linewidth}p{0.18\linewidth}p{0.18\linewidth}@{}}
\toprule
实验 & val last step & best task & last task/success/text & step 200 task/success/text \\
\midrule
GiGPO seed 2026 & 250 & 0.905 & 0.903 / 0.812 / 7.963 & 0.858 / 0.719 / 7.286 \\
GiGPO seed 2077 & 250 & 0.912 & 0.883 / 0.754 / 7.341 & 0.841 / 0.719 / 6.314 \\
GiGPO seed 2501 & 250 & 0.897 & 0.869 / 0.691 / 6.952 & 0.812 / 0.668 / 6.362 \\
StepPPO-v1 seed 2026 & 185 & 0.766 & 0.703 / 0.492 / 3.866 & NA \\
StepPPO-v2 seed 2026 & 250 & 0.806 & 0.662 / 0.547 / 3.485 & 0.562 / 0.469 / 3.000 \\
GiGPO-v1 value-aux detach=false seed 2026 & 200 & 0.802 & 0.545 / 0.398 / 2.950 & 0.545 / 0.398 / 2.950 \\
\bottomrule
\end{longtable}
\end{small}


结论：当前 value-head 相关实验还没有证明收益。StepPPO-v2 的 best task score 到过 0.806，但最终回落到 0.662；GiGPO-value-aux best task score 到过 0.802，但在 step 180--185 出现严重 validation collapse，step 200 仍明显低于原 GiGPO。

\section*{3. Value 收敛情况}

\subsection*{3.1 总体判断}

当前 value head 还不能认为已经稳定收敛。原因有三点：

\begin{enumerate}
\item \texttt{actor/value\_loss} 没有形成单调下降或低方差平台；多个实验中后期仍在 3--5 左右波动。
\item \texttt{actor/value\_pred\_mean} 和 \texttt{actor/value\_return\_mean} 的均值差距已经不大，但这只说明均值校准有所改善，不代表 per-sample value regression 收敛；value loss 仍然较高，说明 target 方差或 outlier 仍很强。
\item value loss 与 policy/format 指标存在同步异常：StepPPO-v2 的 value loss 后期下降，但 valid action ratio 长期偏低；GiGPO-value-aux 保持较高 valid action ratio，却在 step 178--185 出现 response clipping 和 validation 崩盘。
\end{enumerate}

\subsection*{3.2 分段 value 指标}

\begin{small}
\begin{longtable}{@{}p{0.15\linewidth}p{0.15\linewidth}p{0.15\linewidth}p{0.15\linewidth}p{0.15\linewidth}p{0.15\linewidth}@{}}
\toprule
实验 & step 段 & value\_loss & abs(pred-return mean) & value\_clipfrac & valid\_action\_ratio \\
\midrule
StepPPO-v1 & 1--50 & 1.816 & 0.226 & 0.054 & 0.671 \\
StepPPO-v1 & 51--100 & 3.568 & 0.291 & 0.104 & 0.780 \\
StepPPO-v1 & 101--150 & 4.948 & 0.376 & 0.103 & 0.872 \\
StepPPO-v1 & 151--187 & 4.276 & 0.290 & 0.110 & 0.918 \\
StepPPO-v2 & 1--50 & 2.988 & 0.416 & 0.033 & 0.313 \\
StepPPO-v2 & 51--100 & 3.918 & 0.397 & 0.117 & 0.156 \\
StepPPO-v2 & 101--150 & 3.773 & 0.315 & 0.100 & 0.337 \\
StepPPO-v2 & 151--200 & 4.254 & 0.375 & 0.138 & 0.277 \\
StepPPO-v2 & 201--250 & 2.797 & 0.283 & 0.086 & 0.328 \\
GiGPO-value-aux detach=false & 1--50 & 5.098 & 0.484 & 0.078 & 0.917 \\
GiGPO-value-aux detach=false & 51--100 & 6.246 & 0.465 & 0.161 & 0.921 \\
GiGPO-value-aux detach=false & 101--150 & 4.831 & 0.444 & 0.129 & 0.929 \\
GiGPO-value-aux detach=false & 151--200 & 4.573 & 0.399 & 0.126 & 0.877 \\
GiGPO-value-aux detach=false & 201--204 & 4.426 & 0.507 & 0.106 & 0.883 \\
\bottomrule
\end{longtable}
\end{small}


\subsection*{3.3 Last-10 value 状态}

\begin{small}
\begin{longtable}{@{}p{0.11\linewidth}p{0.11\linewidth}p{0.11\linewidth}p{0.11\linewidth}p{0.11\linewidth}p{0.11\linewidth}p{0.11\linewidth}p{0.11\linewidth}@{}}
\toprule
实验 & last-10 step & value\_loss & abs(pred-return mean) & value\_clipfrac & value loss 实际项 & pg\_loss & valid\_action\_ratio \\
\midrule
StepPPO-v1 & 178--187 & 3.552 & 0.243 & 0.098 & 0.355 & -0.005 & 0.935 \\
StepPPO-v2 & 241--250 & 3.011 & 0.380 & 0.093 & 0.090 & -0.085 & 0.281 \\
GiGPO-value-aux detach=false & 195--204 & 4.207 & 0.336 & 0.089 & 0.126 & -0.254 & 0.908 \\
\bottomrule
\end{longtable}
\end{small}


这里的 \texttt{value loss 实际项} 是 \texttt{actor/value\_loss * actor/value\_loss\_coef} 的近似均值。v1 因为 \texttt{value\_loss\_coef=0.1}，后期 value 项约 0.355，数值上远大于 policy loss 均值；v2 和 value-aux 降到 \texttt{0.03} 后，value 项降到 0.09--0.13，但仍可能通过 shared backbone 影响 policy representation。

\section*{4. 关键诊断}

\subsection*{4.1 StepPPO-v1：value loss 对 backbone 干扰偏强}

v1 的 value loss 从 early stage 的 1.816 上升到 101--150 的 4.948，后期仍在 4 左右。由于 \texttt{value\_loss\_coef=0.1} 且 \texttt{detach\_value\_backbone=False}，value regression 会直接更新 actor backbone。结合 v1 validation 明显弱于 GiGPO，可以认为 v1 的 value 学习既没有稳定收敛，也可能干扰了 policy 的 action-format 能力。

\subsection*{4.2 StepPPO-v2：value loss 后期有所缓和，但 policy 行为仍异常}

v2 把 value 系数降到 0.03，并对齐 GiGPO 的 episode advantage 形式。value\_loss 在 201--250 段降到 2.797，比 151--200 段的 4.254 好；但 valid action ratio 长期只有 0.15--0.33，远低于 GiGPO 和 value-aux 的 0.9 左右。因此 v2 的主要瓶颈不只是 value 是否拟合，而是 final advantage / shared-backbone update 后造成了 policy action format 退化。

\subsection*{4.3 GiGPO-value-aux：policy advantage 不变仍会退化}

GiGPO-value-aux 的设计是保持 GiGPO policy advantage 不变，仅新增 value auxiliary target。按理说它应当接近 GiGPO baseline；但实际 step 200 的 task/success/text 只有 0.545/0.398/2.950，而原 GiGPO seed 2026 step 200 是 0.858/0.719/7.286。

这说明退化可以仅由 shared value head 的辅助损失引入，不一定来自 StepPPO 的 mixed advantage。尤其在 step 178--185 附近，GiGPO-value-aux 出现明显异常：

\begin{small}
\begin{longtable}{@{}p{0.15\linewidth}p{0.15\linewidth}p{0.15\linewidth}p{0.15\linewidth}p{0.15\linewidth}p{0.15\linewidth}@{}}
\toprule
step & value\_loss & value\_clipfrac & valid\_action\_ratio & response\_clip\_ratio & validation 附近表现 \\
\midrule
178 & 8.230 & 0.253 & 0.700 & 0.307 & 即将崩盘 \\
179 & 11.031 & 0.345 & 0.653 & 0.352 & 即将崩盘 \\
180 & 6.752 & 0.125 & 0.746 & 0.252 & val task/success/text = 0.153/0.074/0.560 \\
185 & 6.169 & 0.202 & 0.778 & 0.216 & val task/success/text = 0.155/0.102/0.798 \\
\bottomrule
\end{longtable}
\end{small}


该现象提示 value target 的高方差或 clipped value update 的大幅波动，会通过 shared backbone 影响生成分布；即使 valid action ratio 很快恢复，validation 也可能出现短时严重退化。

\section*{5. 当前结论}

\begin{enumerate}
\item GiGPO baseline 已经稳定，仍是当前最强对照。
\item StepPPO-v1 的 value head 没有收敛到可靠 critic，且 value loss coefficient 偏大，干扰风险最高。
\item StepPPO-v2 降低了 value loss 的有效权重，后期 value\_loss 有改善，但 action-format 学习明显退化，最终效果仍弱于 GiGPO。
\item GiGPO-value-aux detach=false 证明：只加 auxiliary value loss，也足以显著破坏 GiGPO；因此当前首要问题是 shared value head / target / backbone gradient 的稳定性，而不是只调 StepPPO advantage mix。
\item 仅看 \texttt{value\_pred\_mean} 与 \texttt{value\_return\_mean} 的均值接近会过于乐观；需要看 RMSE、explained variance、target 分布和 per-step / per-validity 分层误差。
\end{enumerate}

\section*{6. 建议下一步}

\subsection*{6.1 必跑 ablation}

\begin{enumerate}
\item 补跑 \texttt{GiGPO-v1 value-aux detach=true} 的 full run。若 detach=true 明显接近原 GiGPO，则可以确认主要问题来自 value loss 回传 backbone。
\item 补齐正式 \texttt{GRPO seed=2026} baseline，以便区分 non-step baseline、GiGPO step-aware baseline、StepPPO/value-aux 的差异。
\item 对 StepPPO-v2 增加 \texttt{detach\_value\_backbone=True} 或更低 \texttt{value\_loss\_coef=0.01} 的 ablation；这不一定作为最终方法，但可以定位干扰来源。
\end{enumerate}

\subsection*{6.2 必加 value diagnostics}

建议在训练日志中新增以下 value 诊断：

\begin{itemize}
\item \texttt{actor/value\_rmse}、\texttt{actor/value\_mae}、\texttt{actor/value\_explained\_variance}。
\item \texttt{actor/value\_pred\_std}、\texttt{actor/value\_return\_std}、\texttt{actor/value\_return\_p05/p50/p95}。
\item 按 \texttt{step\_idx} 分桶的 value loss / return mean / return std。
\item 按 \texttt{is\_action\_valid} 分层的 value loss，确认 invalid action penalty 是否造成 target outlier。
\item value loss 与 \texttt{response\_length/clip\_ratio}、\texttt{episode/valid\_action\_ratio} 的 rolling correlation。
\end{itemize}

\subsection*{6.3 可能的稳定化方向}

\begin{itemize}
\item 对 value target 做 batch 或 group-level normalization，只把 normalized return 用于 value loss；policy advantage 是否归一化另行控制。
\item 降低或 warmup \texttt{value\_loss\_coef}，例如从 0 到 0.01 / 0.03 线性 warmup，避免早期 value target 噪声破坏 backbone。
\item 比较 \texttt{detach\_value\_backbone=True} 与 \texttt{False}，确认最终方法是否真的需要 value loss 更新 LM backbone。
\item 对 GAE target 做 clipping 或 reward normalization，尤其关注 WebShop sparse/invalid penalty 混合后的 return outlier。
\item 在可视化中叠加 value\_loss、response\_clip\_ratio、valid\_action\_ratio 和 validation task score，重点复盘 GiGPO-value-aux step 170--190 的崩盘窗口。
\end{itemize}

\section*{7. 本文使用的解析口径}

\begin{itemize}
\item 训练行：包含 \texttt{actor/value\_loss} 的 \texttt{step:<global\_step> - ...} console metric 行。
\item validation 行：包含 \texttt{val/success\_rate}、\texttt{val/webshop\_task\_score (not success\_rate)} 或 \texttt{val/text/test\_score} 的 metric 行。
\item \texttt{abs(pred-return mean)}：\texttt{abs(actor/value\_pred\_mean - actor/value\_return\_mean)} 的分段均值，只反映均值校准，不等价于 per-sample RMSE。
\item \texttt{value loss 实际项}：\texttt{actor/value\_loss * actor/value\_loss\_coef}，用于粗略观察 value loss 在总 actor loss 中的数值量级。
\end{itemize}

\end{document}
